\documentclass[11pt]{article}
% ===========================================================================
%  Shared preamble for the AP Calculus AB exam-preparation worksheets.
%  Input right after \documentclass. Files beginning with "_" are not compiled
%  on their own. Navy and gold, to match the course's lesson pages.
% ===========================================================================
\usepackage[letterpaper,top=2.3cm,bottom=2.2cm,left=1.7cm,right=1.7cm,headheight=28pt,footskip=30pt]{geometry}
\usepackage{amsmath,amssymb}
\usepackage[T1]{fontenc}
\usepackage[utf8]{inputenc}
\usepackage{lmodern}
\usepackage{xcolor}
\usepackage[most]{tcolorbox}
\usepackage{enumitem}
\usepackage{fancyhdr}
\usepackage{array}
\usepackage{tikz}
\usepackage{pgfplots}
\pgfplotsset{compat=1.18}

\definecolor{navy}{HTML}{161E54}
\definecolor{ink}{HTML}{27357E}
\definecolor{gold}{HTML}{B98A00}
\definecolor{goldbg}{HTML}{FFFDF3}
\definecolor{inkbg}{HTML}{F6F7FD}
\definecolor{rule}{HTML}{C9CFF0}
\definecolor{okgreen}{HTML}{1B7A44}
\definecolor{okbg}{HTML}{F1F8F3}

\pagestyle{fancy}
\fancyhf{}
\fancyhead[L]{\footnotesize Name: \underline{\hspace{4cm}}}
\fancyhead[C]{\textbf{\color{navy}Integration Academy}\\[-2pt]{\scriptsize Calculus AB \textperiodcentered{} AP\textsuperscript{\textregistered} Exam Preparation}}
\fancyfoot[L]{\tiny\color{black!55}\parbox[t]{0.78\textwidth}{AP\textsuperscript{\textregistered} is a registered trademark of the College Board, which was not involved in the production of, and does not endorse, this material. \textcopyright{} Integration Academy.}}
\fancyfoot[R]{\scriptsize\thepage}
\renewcommand{\headrulewidth}{1.2pt}
\renewcommand{\headrule}{\hbox to\headwidth{\color{gold}\leaders\hrule height \headrulewidth\hfill}}
\renewcommand{\footrulewidth}{0pt}
\setlength{\parindent}{0pt}
\setlength{\parskip}{4pt}

% Title block: \wstitle{code}{title}{unit line}
\newcommand{\wstitle}[3]{\begin{center}{\small\bfseries\color{gold}WORKSHEET #1}\\[3pt]{\LARGE\bfseries\color{navy}#2}\\[3pt]{\small\color{black!60}#3}\end{center}\vspace{2pt}}
% Section heading
\newcommand{\sech}[1]{\par\vspace{10pt}{\large\bfseries\color{navy}#1}\par\nobreak\vspace{2pt}{\color{rule}\hrule height0.8pt}\vspace{7pt}}
% A level tag printed beside a group of questions
\newcommand{\level}[1]{\par\vspace{6pt}{\footnotesize\bfseries\color{gold}\MakeUppercase{#1}}\par\nobreak\vspace{2pt}}

\newtcolorbox{ideabox}{colframe=ink,colback=inkbg,coltitle=white,colbacktitle=ink,fonttitle=\bfseries\small,title={Key ideas},arc=2pt,boxrule=0.8pt,left=8pt,right=8pt}
\newtcolorbox{example}[1]{colframe=rule,colback=white,coltitle=navy,colbacktitle=inkbg,fonttitle=\bfseries\small,title={#1},arc=2pt,boxrule=0.8pt,borderline west={3pt}{0pt}{ink},left=10pt,right=8pt}
\newtcolorbox{frbox}[1]{colframe=gold,colback=goldbg,coltitle=white,colbacktitle=gold,fonttitle=\bfseries\small,title={#1},arc=2pt,boxrule=0.8pt,breakable,left=8pt,right=8pt}
\newtcolorbox{solbox}{colframe=okgreen,colback=okbg,coltitle=white,colbacktitle=okgreen,fonttitle=\bfseries\small,title={Worked solutions},arc=2pt,boxrule=0.8pt,breakable,left=8pt,right=8pt}
\newtcolorbox{scorebox}[1]{colframe=navy,colback=white,coltitle=white,colbacktitle=navy,fonttitle=\bfseries\small,title={#1},arc=2pt,boxrule=0.8pt,breakable,left=8pt,right=8pt}

\newcommand{\soln}{\par\smallskip\textbf{\color{ink}Solution.}\ }
\newcommand{\result}[1]{\par\smallskip\textbf{#1}}
% blank working space of a given height
\newcommand{\work}[1]{\par\vspace{#1}}
% a point value in the margin of a scoring guide
\newcommand{\pts}[1]{\hfill{\footnotesize\bfseries\color{gold}[#1]}}

% Coordinate grid for graph questions: \begin{apgraph}{xmin}{xmax}{ymin}{ymax} ... \end{apgraph}
\newenvironment{apgraph}[4]{\begin{center}\begin{tikzpicture}\begin{axis}[
  axis lines=middle, xmin=#1, xmax=#2, ymin=#3, ymax=#4,
  grid=both, grid style={black!12}, xtick distance=1, ytick distance=1,
  ticklabel style={font=\scriptsize}, xlabel={$x$}, ylabel={$y$},
  width=10cm, height=7.2cm, samples=80, clip=false,
  every axis plot/.append style={line width=1.1pt, color=ink}]}{\end{axis}\end{tikzpicture}\end{center}}
\newcommand{\closeddot}[2]{\addplot[only marks, mark=*, mark size=2.2pt, color=ink] coordinates {(#1,#2)};}
\newcommand{\opendot}[2]{\addplot[only marks, mark=*, mark size=2.2pt, color=ink, mark options={fill=white}] coordinates {(#1,#2)};}

% Multiple-choice options: four across, or stacked when the options are long.
\newcommand{\opts}[4]{\par\vspace{3pt}\begin{tabular}{@{}p{0.235\linewidth}p{0.235\linewidth}p{0.235\linewidth}p{0.235\linewidth}@{}}(A)\ #1 & (B)\ #2 & (C)\ #3 & (D)\ #4\end{tabular}\par\vspace{9pt}}
\newcommand{\optsv}[4]{\par\vspace{3pt}\begin{tabular}{@{}l@{\ \ }p{0.9\linewidth}@{}}(A) & #1\\(B) & #2\\(C) & #3\\(D) & #4\end{tabular}\par\vspace{9pt}}

\fancyhead[R]{\footnotesize\color{black!60}Worksheet 0.3}
\begin{document}

\wstitle{0.3}{Exponentials, Logarithms and Their Laws}{Unit 0: Readiness}

\begin{ideabox}
\begin{itemize}[leftmargin=*,itemsep=2pt,topsep=2pt]
  \item $y=\ln x$ means $e^{y}=x$. So $\ln\!\left(e^{x}\right)=x$ for all $x$, and $e^{\ln x}=x$ for $x>0$.
  \item $\ln(ab)=\ln a+\ln b$, \quad $\ln\dfrac ab=\ln a-\ln b$, \quad $\ln\!\left(a^{b}\right)=b\ln a$. There is no law for $\ln(a+b)$.
  \item To solve for an exponent: \textbf{isolate} the exponential, then take the natural logarithm.
  \item A logarithmic equation can produce a value that is not allowed. \textbf{Check} every solution in the original equation.
\end{itemize}
\end{ideabox}

\sech{Worked examples}
\begin{example}{Example 1 \textperiodcentered{} Simplifying}
Simplify $e^{\ln4+\ln3}$ and $\ln\!\left(e^{2}\sqrt{e}\right)$.\soln $e^{\ln4+\ln3}=e^{\ln12}$. $\ln\!\left(e^{2}\cdot e^{1/2}\right)=\ln\!\left(e^{5/2}\right)$.\result{$12$ and $\dfrac52$.}
\end{example}

\begin{example}{Example 2 \textperiodcentered{} Isolate, then take logs}
Solve $5e^{2x}-3=17$.\soln $5e^{2x}=20$, so $e^{2x}=4$ and $2x=\ln4$.\result{$x=\dfrac{\ln4}{2}=\ln2$.}
\end{example}

\begin{example}{Example 3 \textperiodcentered{} A difference of logarithms}
Solve $\ln(x+6)-\ln x=\ln4$.\soln $\ln\dfrac{x+6}{x}=\ln4$, so $x+6=4x$.\result{$x=2$, which makes both logarithms defined.}
\end{example}

\newpage
\sech{Practice}

\level{Core}
\begin{enumerate}[leftmargin=*,itemsep=2pt]
  \item Simplify $\ln\!\left(e^{7}\right)-e^{\ln2}$.\work{2.3cm}
  \item Solve $e^{x}=\dfrac12$.\work{2.3cm}
  \item Solve $\ln x=-1$.\work{2.3cm}
  \item Write $2\ln a-\ln b+\ln c$ as a single logarithm.\work{2.3cm}
\end{enumerate}

\level{Exam level}
\begin{enumerate}[leftmargin=*,itemsep=2pt,start=5]
  \item Solve $e^{2x}-9e^{x}+20=0$.\work{2.8cm}
  \item Solve $\ln(x-2)+\ln(x+2)=\ln5$.\work{2.8cm}
  \item Solve $3^{x}=15$, giving an exact answer in terms of natural logarithms.\work{2.8cm}
  \item Write $\ln\dfrac{x^3e^{x}}{\sqrt{x^2+1}}$ as a sum and difference of simpler terms.\work{2.8cm}
\end{enumerate}

\begin{frbox}{Question 9 \textperiodcentered{} Free-response practice}
The amount of a drug in a patient's bloodstream $t$ hours after an injection is $A(t)=80e^{-0.1t}$ milligrams.
\begin{enumerate}[label=(\alph*),leftmargin=*,itemsep=2pt]
  \item How much of the drug is present at $t=0$?\work{2.6cm}
  \item Find the exact time at which 20 milligrams remain.\work{2.6cm}
  \item Find the exact half-life of the drug, and explain how it is related to your answer to part (b).\work{2.6cm}
\end{enumerate}
\end{frbox}

\newpage
\begin{solbox}
\begin{enumerate}[leftmargin=*,itemsep=4pt]
  \item $7-2=5$.
  \item $x=\ln\dfrac12=-\ln2$.
  \item $x=e^{-1}=\dfrac1e$.
  \item $\ln\dfrac{a^2c}{b}$.
  \item $(e^{x}-4)(e^{x}-5)=0$, so $x=\ln4$ or $x=\ln5$.
  \item $\ln\!\left(x^2-4\right)=\ln5$ gives $x^2=9$. Only $x=3$ is allowed, since $x=-3$ makes the logarithms undefined.
  \item $x\ln3=\ln15$, so $x=\dfrac{\ln15}{\ln3}$.
  \item $3\ln x+x-\dfrac12\ln\!\left(x^2+1\right)$.
  \item \begin{enumerate}[label=(\alph*),leftmargin=*,itemsep=3pt]
    \item $A(0)=80$ milligrams.
    \item $80e^{-0.1t}=20$ gives $e^{-0.1t}=\dfrac14$, so $t=10\ln4$ hours (about 13.9).
    \item $e^{-0.1t}=\dfrac12$ gives $t=10\ln2$ hours. Falling from 80 to 20 is two halvings, and $10\ln4=2\cdot10\ln2$.
  \end{enumerate}
\end{enumerate}
\end{solbox}

\end{document}
